% Source: https://artofproblemsolving.com/community/c3h1688488_help_pls
% Stable identity: aops-topic-1688488-q01
Work from left to right: $-2-(-1)-3=-2+1-3=\boxed{-4}$.

The student's expression is $-2-\bigl(-1+(-3)\bigr)=-2-(-4)=2$, which is different. Subtracting a bracket subtracts \emph{every} term inside it: $-2-(-1+(-3))=-2+1+3$. In the original only $-1$ is subtracted, while the $3$ is subtracted separately; grouping $-1$ and $-3$ inside one bracket after a minus sign turns the subtraction of $3$ into an addition of $3$. A correct grouping is
\[
-2-(-1)-3=-2-\bigl((-1)+3\bigr)=-2-2=-4 ,
\]
since subtracting $-1$ and then $3$ is the same as subtracting their sum $(-1)+3=2$.
